% ==============================================================================
\chapter{Implementación}
\label{ch:implementacion}
% ==============================================================================

Este capítulo documenta el desarrollo. Describe la estructura del repositorio, el entorno, los seis sprints (con las decisiones técnicas más relevantes de cada uno), algunos extractos de código y el \textit{pipeline} de integración continua.


\section{Entorno y herramientas}
\label{sec:entorno}

He trabajado en Windows~11 con PowerShell. Las herramientas usadas son:

\begin{itemize}[leftmargin=1.2cm, itemsep=2pt]
  \item \textbf{IDE}: Google Antigravity (\textit{fork} de VS~Code con agentes de IA).
  \item \textbf{Git} con convención \textit{Conventional Commits}.
  \item \textbf{Docker Desktop} con Docker Compose v2.
  \item \textbf{Node.js 20 LTS} y \textbf{PostgreSQL 16 + PostGIS 3.4} (en contenedor).
  \item \textbf{Pruebas}: Jest~29 + Supertest~7 (\textit{backend}); Playwright~1.38 (E2E, vía MCP de Antigravity).
  \item \textbf{Bibliografía}: Zotero exportando a BibLaTeX.
  \item \textbf{Documentación API}: Swagger UI en \texttt{/api/docs}.
  \item \textbf{MCPs}: GitHub, Context7, Playwright, Figma.
\end{itemize}

Elegí Antigravity por la integración nativa con agentes (que encaja con el flujo SDD descrito en la sección~\ref{sec:metodologia}) y la compatibilidad con el ecosistema MCP, que permite ejecutar Playwright sin salir del editor.


\section{Estructura del repositorio}
\label{sec:estructura-repo}

\begin{lstlisting}[basicstyle=\ttfamily\footnotesize, frame=single, backgroundcolor=\color{gray!5}]
tfg-medioambiente-ua/
|-- .github/workflows/       # CI/CD GitHub Actions
|-- database/                # init.sql, seeds
|-- docs/                    # Documentacion del TFG
|   |-- memoria/             #   Fuentes LaTeX
|   |-- diagramas/           #   Fuentes PlantUML
|   |-- sprints/             #   Specs SDD por sprint
|-- nginx/                   # Reverse proxy
|-- src/
|   |-- backend/             # API Node.js + Express
|   `-- frontend/            # SPA React + Tailwind + PWA
|-- docker-compose.dev.yml
|-- docker-compose.yml
|-- Makefile
`-- EcoAlerta_Development_Spec_v1.md
\end{lstlisting}

El \textit{backend} sigue la organización \textit{routes~$\rightarrow$~controllers~$\rightarrow$~services} típica de Express: las rutas definen \textit{endpoints}, los controladores orquestan petición/respuesta, y los servicios encapsulan la lógica y las consultas. Los validadores usan \texttt{express-validator}. El \textit{frontend} se organiza por \textit{feature}: \path{components/}, \path{pages/}, \path{context/}, \path{hooks/}, \path{services/}, \path{utils/}.


\section{Desarrollo por sprints}
\label{sec:sprints}

El desarrollo se organizó en seis \textit{sprints} de una semana. La Tabla~\ref{tab:sprints-resumen} resume cada uno; abajo amplío los puntos técnicos más relevantes.

\begin{table}[H]
\centering
\caption{Resumen de los seis sprints.}
\label{tab:sprints-resumen}
\small
\begin{tabular}{@{}L{1cm}L{3.8cm}L{8cm}@{}}
\toprule
\textbf{S.} & \textbf{Tema} & \textbf{Entregables} \\
\midrule
1 & Auth + BD              & Esquema PostGIS, JWT dual, roles, \textit{middlewares} de seguridad \\
2 & CRUD incidencias       & Creación con foto y GPS, validaciones, \textit{thumbnails} con Sharp \\
3 & Mapa interactivo       & Leaflet + \textit{clusters}, filtros, búsqueda con \texttt{ST\_DWithin} \\
4 & Admin + delegación     & Dashboard, cambio de estado, asignación, notificaciones \\
5 & Testing + PWA          & Jest, Playwright, Service Worker, manifest, Swagger \\
6 & \textit{Security hardening} & 2FA TOTP, Turnstile, \textit{audit log}, \textit{rate limit} por usuario \\
\bottomrule
\end{tabular}
\end{table}


\subsection{Sprint 1: autenticación y BD}

La BD se inicializa con un \texttt{init.sql} ejecutado por el contenedor en el primer arranque: habilita PostGIS, crea las trece tablas, define los \texttt{ENUM} (\texttt{user\_role}, \texttt{incident\_status}, \texttt{incident\_severity}), crea el índice GIST sobre \texttt{incidents.location} e inserta las 12 categorías y las 5 entidades por defecto.

Para autenticación uso JWT con \textit{access token} de 15~minutos y \textit{refresh token} de 7~días. La función que firma ambos:

\begin{lstlisting}[language=JavaScript, basicstyle=\ttfamily\footnotesize, frame=single, backgroundcolor=\color{gray!5}]
const generateTokens = (user) => {
  const payload = { id: user.id, email: user.email, role: user.role };
  return {
    accessToken:  jwt.sign(payload, process.env.JWT_SECRET, { expiresIn: '15m' }),
    refreshToken: jwt.sign(payload, process.env.JWT_SECRET, { expiresIn: '7d' }),
  };
};
\end{lstlisting}

Las contraseñas se guardan en \texttt{password\_hash} con \textit{bcrypt} a 10 rondas (RNF-05). El \textit{middleware} \texttt{authenticate} verifica el token en cada petición protegida y pone el \textit{payload} en \texttt{req.user}; \texttt{optionalAuth} permite que los \textit{endpoints} de lectura pública funcionen sin token (RF-03). El control por rol es un \textit{factory}:

\begin{lstlisting}[language=JavaScript, basicstyle=\ttfamily\footnotesize, frame=single, backgroundcolor=\color{gray!5}]
const requireRole = (roles) => (req, res, next) => {
  if (!req.user || !roles.includes(req.user.role))
    return res.status(403).json({ success: false, error: 'Acceso restringido' });
  next();
};
const requireAdmin  = requireRole(['admin']);
const requireEntity = requireRole(['entity']);
\end{lstlisting}


\subsection{Sprint 2: CRUD de incidencias}

Aquí se monta el flujo principal: crear una incidencia con foto y GPS. Tres puntos a destacar:

\begin{itemize}[leftmargin=1.2cm, itemsep=2pt]
  \item \textbf{Subida \textit{multipart}} con \texttt{multer}, controlando tamaño y tipo MIME antes del controlador.
  \item \textbf{\textit{Thumbnails}} JPEG (400~px, calidad 80) con \texttt{sharp}, generados en el servidor (RF-16).
  \item \textbf{Ubicación} guardada con \texttt{ST\_SetSRID(ST\_MakePoint(lng, lat), 4326)} para fijar el SRID y evitar ambigüedad en consultas posteriores.
\end{itemize}

En el \textit{frontend}, la página \texttt{CreateIncidentPage} es un formulario por pasos: geolocalización automática (con \textit{fallback} a selección manual sobre el mapa), categoría y severidad, fotos con previsualización, descripción. Si no hay conexión, el borrador queda en IndexedDB y se sincroniza cuando vuelve la red.


\subsection{Sprint 3: mapa interactivo}

Uso \texttt{react-leaflet} sobre \textit{tiles} de OpenStreetMap. Los marcadores se colorean por severidad mediante iconos SVG (RF-23) y se agrupan en \textit{clusters} con \texttt{leaflet.markercluster} (RF-21). Los filtros se componen dinámicamente y se proyectan en la consulta al \textit{backend}.

La búsqueda por radio es la consulta más característica del sistema:

\begin{lstlisting}[language=JavaScript, basicstyle=\ttfamily\footnotesize, frame=single, backgroundcolor=\color{gray!5}]
const getNearbyIncidents = async (lat, lng, radiusKm = 5, filters = {}) => {
  const radius = radiusKm * 1000;
  const where = [`ST_DWithin(i.location::geography,
                  ST_SetSRID(ST_MakePoint($1, $2), 4326)::geography, $3)`];
  const values = [lng, lat, radius];
  // ... filtros dinamicos por status, categoryId, severity ...
  const result = await query(`SELECT i.*, ST_Distance(...) AS distance_meters
    FROM incidents i WHERE ${where.join(' AND ')}
    ORDER BY distance_meters ASC LIMIT 50`, values);
  return result.rows;
};
\end{lstlisting}

Tres detalles importantes: la proyección a \texttt{geography} hace que la distancia sea en metros sobre la superficie terrestre y no en grados; los votos se agregan en una subconsulta para no contar filas duplicadas; el orden por \texttt{distance\_meters} aprovecha que la distancia ya se calculó. El \textit{geocoding} por dirección (RF-24) llama a Nominatim con caché en memoria.


\subsection{Sprint 4: admin y delegación}

Páginas \texttt{AdminDashboardPage} (con Recharts), \texttt{AdminIncidentsPage} (tabla con filtros y acciones masivas), \texttt{AdminUsersPage} y \texttt{AdminEntitiesPage}. La vista para entidades responsables es \texttt{EntityDashboardPage}: filtra automáticamente por \texttt{assigned\_entity\_id} del usuario.

El cambio de estado (RF-31) es una transición atómica: la función \texttt{changeIncidentStatus} valida que la transición sea coherente con el diagrama de la Figura~\ref{fig:estados-incidencia}, escribe en \texttt{incident\_status\_history}, actualiza el campo \texttt{status} y dispara una notificación al autor y a los seguidores. Todo dentro de \texttt{BEGIN}/\texttt{COMMIT} para asegurar consistencia.

Las notificaciones se generan en el \textit{backend} y se entregan por \textit{polling} cada 30 segundos desde el componente \texttt{NotificationBell}. Elegí \textit{polling} sobre WebSockets por simplicidad de despliegue y porque la frecuencia es baja. Las notificaciones \textit{push} reales (RF-45) se montaron en el siguiente sprint con el Service Worker.


\subsection{Sprint 5: testing y PWA}

Aquí consolidé la calidad. Bateria de pruebas en tres niveles (unitarias con Jest, integración con Supertest, E2E con Playwright). La cobertura final del \textit{backend} se discute en el Capítulo~\ref{ch:evaluacion}. Quedó por debajo del 70\,\%, sobre todo en \textit{branches}, y lo recojo como limitación.

PWA: el \texttt{manifest.json} declara iconos, nombre, color y \texttt{display: standalone}. El Service Worker usa \textit{cache-first} para \textit{assets} estáticos y \textit{network-first} con \textit{fallback} a caché para llamadas API. Esto permite consultar el mapa y las incidencias guardadas sin conexión.

Swagger se genera con \texttt{swagger-jsdoc} a partir de anotaciones \texttt{@swagger} en las rutas y se sirve en \texttt{/api/docs} con \texttt{swagger-ui-express}. Documentación y código no se desincronizan.


\subsection{Sprint 6: \textit{security hardening}}

Aproveché el tiempo extra para añadir tres mejoras de seguridad. Las tablas \texttt{user\_2fa}, \texttt{user\_recovery\_codes} y \texttt{security\_audit\_log} entraron aquí, pasando el modelo de 10 a 13 entidades.

\begin{itemize}[leftmargin=1.2cm, itemsep=2pt]
  \item \textbf{2FA TOTP}. Integración de \texttt{otpauth} (RFC~6238) y \texttt{qrcode}. El secreto se guarda \textbf{cifrado} con AES-256 a partir de una clave maestra fuera del código fuente. Los códigos de recuperación, de un solo uso, se guardan \textit{hasheados}.
  \item \textbf{Turnstile} (\textit{captcha} de Cloudflare). Componente \texttt{@marsidev/react-turnstile} en el \textit{frontend}; \textit{middleware} \texttt{turnstile.middleware.js} valida el token contra la API de Cloudflare en el \textit{backend}.
  \item \textbf{\textit{Audit log}}. Cada evento de seguridad (login OK/KO, cambio de rol, activación 2FA) se guarda en \texttt{security\_audit\_log} con IP, \textit{user agent} y metadatos JSONB. No se borra. \texttt{audit.service.js} centraliza la escritura.
\end{itemize}

Añadí también \textit{rate limiting} por usuario con los contadores \texttt{failed\_login\_attempts} y \texttt{failed\_login\_window\_start} en \texttt{users}, complementando el de \texttt{express-rate-limit} por IP. La doble protección frena fuerza bruta aunque el atacante rote IPs.


\section{\textit{Frontend}: notas adicionales}
\label{sec:frontend-destacados}

La aplicación usa \texttt{react-router-dom} v6 con rutas públicas (mapa, detalle, login, registro) y rutas protegidas por rol. Las protegidas se envuelven en un \texttt{ProtectedRoute} que mira el contexto y redirige a \texttt{/login} si toca. \textit{Layout} con \texttt{Navbar} en escritorio y \texttt{MobileBottomNav} en móvil.

El estado de autenticación vive en \texttt{AuthContext}, que expone \texttt{login}, \texttt{logout}, \texttt{refresh} y \texttt{user}. Los \textit{tokens} van en \texttt{localStorage}. Un interceptor de Axios detecta 401 por expiración y renueva el \textit{access token} con el \textit{refresh} antes de reintentar la petición. El \textit{hook} \texttt{useAuth} es el punto de acceso al contexto.

Componentes reutilizables clave: \texttt{SeverityBadge}, \texttt{StatusBadge}, \texttt{IncidentMarker}, \texttt{MapFilters}, \texttt{NotificationBell}, \texttt{TwoFactorSetupModal}, \texttt{TwoFactorVerifyForm}.


\section{Despliegue con Docker Compose}
\label{sec:despliegue}

\path{docker-compose.yml} define cinco componentes: \texttt{db} (PostgreSQL 16 + PostGIS 3.4 con \texttt{init.sql} y volumen \texttt{pgdata}), \texttt{backend} (Node 20 con Express en el puerto 5000), \texttt{frontend} (build de la SPA), \texttt{nginx} (\textit{reverse proxy} HTTPS y servidor de estáticos), y el volumen \texttt{pgdata}. \texttt{docker compose up -d} levanta el sistema completo (RNF-10).

El \texttt{Makefile} ofrece atajos: \texttt{make dev}, \texttt{make test}, \texttt{make lint}, \texttt{make prod}, \texttt{make db-reset}.


\section{Integración continua}
\label{sec:cicd}

\path{.github/workflows/ci.yml} se dispara en cada \textit{push} y \textit{pull request}. Tres etapas:

\begin{enumerate}[leftmargin=1.2cm, itemsep=2pt]
  \item \textbf{Lint} con ESLint sobre \textit{backend} y \textit{frontend}.
  \item \textbf{Test} con Jest del \textit{backend} y reporte de cobertura archivado como artefacto.
  \item \textbf{Build} de la imagen de producción del \textit{frontend}.
\end{enumerate}

El historial de \textit{commits} sigue \textit{Conventional Commits}, lo que permite generar el \texttt{CHANGELOG.md} automáticamente.


\section{Capturas de pantalla}
\label{sec:sprints-ui}

% TODO: incorporar capturas reales en figuras/ui-*.png

\begin{figure}[H]
  \centering
  \begin{subfigure}[b]{0.48\textwidth}\centering
    \fbox{\parbox[b][4cm][c]{\linewidth}{\centering\textit{ui-mapa-principal.png}}}
    \caption{Mapa principal.} \label{fig:ui-mapa}
  \end{subfigure}
  \hfill
  \begin{subfigure}[b]{0.48\textwidth}\centering
    \fbox{\parbox[b][4cm][c]{\linewidth}{\centering\textit{ui-crear-incidencia.png}}}
    \caption{Formulario de reporte.} \label{fig:ui-reporte}
  \end{subfigure}
  \caption{Vistas del ciudadano. (Capturas pendientes.)}
\end{figure}


\section{Retos técnicos resueltos}
\label{sec:retos}

Algunos retos no triviales del proyecto, relevantes para la defensa:

\paragraph{Rendimiento geoespacial.}  La consulta «incidencias en un radio» se degrada rápido sin índice. \texttt{ST\_DWithin} con índice GIST hace que PostgreSQL descarte filas antes de calcular distancias, manteniendo latencias por debajo del umbral del RNF-01 incluso con miles de incidencias.

\paragraph{Renovación transparente del \textit{token}.}  JWT \textit{stateless} simplifica escalar pero complica invalidar. Lo resuelvo con \textit{access} corto (15~min), \textit{refresh} largo (7~días) con rotación en cada uso, y un interceptor Axios que renueva el \textit{access} sin que el usuario vea un 401 espurio.

\paragraph{Inyección SQL.}  Todas las consultas usan parámetros posicionales (\texttt{\$1}, \texttt{\$2}\dots) del \textit{driver} \texttt{pg}, nunca concatenación. En \texttt{getNearbyIncidents}, donde el número de filtros varía, los valores se acumulan en un \textit{array} y se pasan al \textit{driver} sin entrar al \textit{parser} SQL como cadena.

\paragraph{Cifrado de secretos 2FA.}  Un secreto TOTP en claro en BD sería como una contraseña débil ante un \textit{leak}. Lo cifro con AES-256 usando una clave maestra en variable de entorno, así un volcado de BD no basta para comprometer el segundo factor. Los códigos de recuperación van \textit{hasheados} porque nunca tengo que recuperarlos en claro.

\paragraph{\textit{Offline} con sincronización diferida.}  El uso típico ocurre en zonas con cobertura inestable (rural, parajes naturales). El Service Worker cachea la UI y las llamadas a \texttt{/incidents/nearby}; los borradores van a IndexedDB y se sincronizan al recuperar conexión. Activación automática en navegadores con \textit{Background Sync}.


\section{Asistencia de IA en la implementación}
\label{sec:ia-implementacion}

Como anticipé en la sección~\ref{sec:uso-ia}, parte importante del código se escribió con asistencia del agente de Antigravity siguiendo SDD. El flujo en cada sprint fue:

\begin{enumerate}[leftmargin=1.2cm, itemsep=2pt]
  \item Redacto la especificación del sprint en \path{docs/sprints/sprint-N.md}: requisitos a cubrir, criterios de aceptación, esquema de pruebas.
  \item El agente genera código a partir de la especificación, con acceso \textit{read-only} al resto del repositorio como contexto.
  \item Reviso, ejecuto pruebas locales, depuro y ajusto manualmente lo que el agente hace mal.
  \item \textit{Commit} en \textit{Conventional Commits} y \texttt{CHANGELOG.md} actualizado.
\end{enumerate}

Las intervenciones manuales más relevantes ocurrieron en la consulta geoespacial (el agente proponía \texttt{ST\_Distance} con \texttt{WHERE}, ineficiente; lo cambié a \texttt{ST\_DWithin} tras revisar la documentación de PostGIS), en la transacción del cambio de estado (añadí \texttt{BEGIN}/\texttt{COMMIT} que el agente había omitido) y en la configuración del Service Worker (donde la estrategia inicial rompía las peticiones autenticadas). En esos puntos la IA propone soluciones que pasan los tests pero no son las óptimas; ahí entra la revisión humana.
